
\section{Constraints}

In view of \eqref{cap3_x_in_X} and \eqref{cap3_u_in_U}, the constraints on the optimization variables ca be expressed as

\noindent
\begin{align}
	\mathbf{\hat{u}}[k]&= \left[
		\begin{matrix}
			\hat{u}[k | k] \\
			\hat{u}[k + 1 | k] \\
			\vdots \\
			\hat{u}[k + N-1 | k]
		\end{matrix}
	\right]
	\in \mathcal{U} \times  \mathcal{U} \times \cdots \times \mathcal{U} = \mathcal{U}^N \\
		\mathbf{\hat{x}}[k] &= \left[
		\begin{matrix}
			\hat{x}[k + 1 | k] \\
			\hat{x}[k + 2 | k] \\
			\vdots \\
			\hat{x}[k + N | k]
		\end{matrix}
	\right]
	\in \mathcal{X} \times  \mathcal{X} \times \cdots \times \mathcal{X}_f = \mathcal{X}^{N-1} \times \mathcal{X}_f
\end{align}

\noindent
where

\noindent
\begin{equation}
	\mathcal{X}_f = \big\{x \in \mathbb{R}^n : S_f x \le b_f \big\}
\end{equation}
is the maximum invariant and admissible set (MAS) concerning the terminal state and control
constraints under the infinite-horizon DLQR control law.

\noindent
\begin{equation}
	\left\{
		\hat{x} \in \mathbb{R}^n : S_f \hat{x} \leq b_f
	\right\}
\end{equation}

\noindent
The constraint on \highlight{$\mathbf{\hat{u}}$} can be  written as

\noindent
\begin{equation}
	\left[
		\begin{matrix}
			S_u & 0 & \cdots & 0 \\
			0 & S_u & \cdots & 0 \\
			\vdots & \vdots & \ddots & \vdots \\
			0 & 0 & \cdots & S_u
		\end{matrix}
	\right]
	\left[
		\begin{matrix}
			\hat{u}[k | k] \\
			\hat{u}[k + 1 | k] \\
			\vdots \\
			\hat{u}[k + N-1 | k]
		\end{matrix}
	\right]
	\leq
	\left[
		\begin{matrix}
			b_u \\
			b_u \\
			\vdots \\
			b_u
		\end{matrix}
	\right]
\end{equation}

\noindent
or, in a shorter form

\noindent
\begin{equation}
	\mathbf{S_u \hat{u}} \leq \mathbf{b_u}
\end{equation}

\noindent
with

\noindent
\begin{equation}
	\mathbf{S_u} = diag_N(S_u), \mathbf{b_u} = col_N(b_u)
\end{equation}

Similarly, in light of \eqref{cap3_eq_Sx_bx}, the constraints on \highlight{$\mathbf{\hat{x}}$} can be written as

\begin{equation}
	\left[
		\begin{matrix}
			S_x & 0 & \cdots & 0 \\
			0 & S_x & \cdots & 0 \\
			\vdots & \vdots & \ddots & \vdots \\
			0 & 0 & \cdots & S_x
		\end{matrix}
	\right]
	\left[
		\begin{matrix}
			\hat{x}[k + 1 | k] \\
			\hat{x}[k + 2 | k] \\
			\vdots \\
			\hat{x}[k + N | k]
		\end{matrix}
	\right]
	\leq
	\left[
		\begin{matrix}
			b_x \\
			b_x \\
			\vdots \\
			b_x
		\end{matrix}
	\right]
\end{equation}

\noindent
or, in short

\noindent
\begin{equation}
	\mathbf{S_x \hat{x}} \leq \mathbf{b_x}
\end{equation}

\noindent
with

\noindent
\begin{equation}
	\mathbf{S_x} = 
	\left[
		\begin{matrix}
			diag_{N-1}(S_x) & 0 \\
			0 & S_f
		\end{matrix}
	\right], \mathbf{b_x} = 
	\left[
		\begin{matrix}
			col_{N-1}(b_x) \\
			b_f
		\end{matrix}
	\right]
\end{equation}

\comentarioblue{
	essa discussao tem que ser apresentada depois de se deduzir a equacao de predicao
	na secao 3.3 $\downarrow \downarrow$
}

The vectors \highlight{$\mathbf{\hat{x}}[k|k]$} and \highlight{$\mathbf{\hat{u}}[k|k]$} are linked by

\begin{equation}
	\mathbf{\hat{x}}[k+1|k] = A \mathbf{\hat{x}}[k|k] + B \mathbf{\hat{u}}[k|k]
\end{equation}






\noindent
so the constraint

\begin{equation}
	\mathbf{S_x} \mathbf{x} \le \mathbf{b_x}
\end{equation}

\noindent
can be rewritten as 


\begin{equation}
	\mathbf{S_x} \mathbf{B} \mathbf{u} \le \mathbf{b_x} - \mathbf{S_x} \mathbf{A} x(0)
\end{equation}

Combined with \highlight{$\mathbf{S_u} \mathbf{u} \le \mathbf{b_u}$}, it can be expressed in the form

\begin{equation}
	\mathbf{S} \mathbf{u} \le \mathbf{b}
\end{equation}

\noindent
with

\begin{equation}
	\mathbf{S} = \left[ \begin{matrix} \mathbf{S_x} \mathbf{B} \\ \mathbf{S_u} \end{matrix} \right], \mathbf{b} = 
	\left[ \begin{matrix} \mathbf{b_x} - \mathbf{S_x} \mathbf{A} \hat{x}[k|k] \\ \mathbf{b_u} \end{matrix} \right]
\end{equation}

\comentarioblue{
	essa discussao tem que ser apresentada depois de se deduzir a equacao de predicao
	na secao 3.3 $\uparrow \uparrow$
}
